\documentclass[12pt,a4paper]{article}
\usepackage{ra_preamble}
\title{\textbf{Relational Actualism: Irreversible Events as the\\
Primitive Basis of Physical Reality}\\[0.4em]
\large\normalfont Relational Actualism Paper~V — The Complete Framework}
\author{Joshua F.~Sandeman\\[0.3em]
\small Independent Researcher, Salem, Oregon\\
\small \texttt{sansuikyo@gmail.com} $\cdot$ \texttt{relationalactualism.org}}
\date{April 2026}

\begin{document}
\maketitle
\begin{abstract}
Relational Actualism (RA) is a framework in which irreversible on-shell actualization events --- physical interactions producing a permanent increase in quantum relative entropy $\Delta S(\rho\|\sigma_0) > \DeltaS = 0.60069\,\mathrm{nats}$ --- are the primitive physical reality.  Spacetime is the growing directed acyclic graph (DAG) of these events.  From a single ontological commitment, without free parameters, RA derives: $\alpha_{\mathrm{EM}}=1/137$ \statusLV, $\alpha_s=1/\sqrt{72}$ \statusCV, the Koide formula $K=2/3$ \statusLV, the complete Standard Model spectrum in five topology classes \statusLV, $\Lambda=0$ structurally \statusDR, the BMV null prediction ($\tact \approx 1.9\times10^4\,\mathrm{s} \gg \tau_{\mathrm{exp}}$) \statusDR, the Hubble tension ($H_0^{\mathrm{Eridanus}}\approx75.9\,\mathrm{km/s/Mpc}$) \statusCV, the WIMP prohibition \statusAR, the Boltzmann Brain prohibition \statusDR, the minimum viable universe seed ($N_{\min}=2$ connected, $V=+9$) \statusDR, the universal decoherence timescale ($\tau_d=20\,\mathrm{fs}$ at $300\,\mathrm{K}$) \statusDR, the Kinematic Coherence Bound ($N_{\max}=\eta\cdot p_{\mathrm{th}}$) \statusLV, and the Causal Firewall origin-of-life threshold (same $N_{\min}=2$ as the Big Bang) \statusDR.  Approximately $110$ results are formally verified in Lean~4 with zero sorry tags in the core files.  All claims are explicitly marked with epistemic status.
\end{abstract}

\tableofcontents

%───────────────────────────────────────────────────────────────
\section{The Single Commitment}

\subsection{The Ontological Commitment}

\begin{mdframed}[linecolor=cv_blue!60,backgroundcolor=cv_blue!5,linewidth=1.5pt]
\textbf{The Relational Actualism Commitment:} An interaction constitutes a \emph{physical event} if and only if it produces an irreversible increase in quantum relative entropy $\Delta S(\rho\|\sigma_0) > \DeltaS = 0.60069\,\mathrm{nats}$ relative to the vacuum state~$\sigma_0$.  Such events write permanent directed edges into the growing causal DAG.  Everything else is derived.
\end{mdframed}

\subsection{Time and the Arrow}

Time is not a background dimension.  Time is the process of the causal DAG growing.  The arrow of time is not explained by low-entropy initial conditions --- it is encoded in the framework's most primitive commitment, since actualization is by definition irreversible.  Proper time along a worldline is the integer count of actualization events along that path.  Relativistic time dilation is the statement that different worldlines accumulate events at different rates.

\subsection{The BDG Action}

The dynamics are governed by the Benincasa--Dowker--Glaser action in $d=4$:
\begin{equation}
  \SBDG
  \label{eq:BDG}
\end{equation}
where $N_k$ counts causal-set elements at depth $k$ from candidate vertex $v$.  The integers $(1,-1,9,-16,8)$ are uniquely determined by the requirement that this operator converges to the d'Alembertian in the Lorentzian continuum limit \cite{Dowker2013}.  They are \emph{not} free parameters.

%───────────────────────────────────────────────────────────────
\section{The Epistemic Map}

Every claim in RA carries an explicit epistemic status:

\begin{table}[h]
\centering
\caption{Epistemic status categories}
\begin{tabular}{clp{8cm}}
\toprule
Badge & Name & Meaning \\
\midrule
\statusLV & Lean-Verified & Machine-checked in Lean~4, zero sorry tags.  Certainty: as high as mathematics allows. \\
\statusCV & Computation-Verified & Certified numerical computation with explicit error bounds. \\
\statusDR & Derived & All steps explicit in the papers; argument complete but not machine-checked. \\
\statusAR & Argued & Key steps identified; some details incomplete.  Upgradeable. \\
\statusOP & Open & Precisely scoped gap with stated path to resolution. \\
\bottomrule
\end{tabular}
\end{table}

This transparency is intentional.  Distinguishing what is \emph{proved} from what is \emph{argued} builds credibility and identifies exactly where theoretical effort should be directed.

%───────────────────────────────────────────────────────────────
\section{Lean~4 Formal Verification}

\subsection{Core Results (Zero Sorry Tags)}

The following are machine-verified in RA\_D1\_Proofs.lean, RA\_AmpLocality.lean, RA\_Spin2\_Macro.lean, and RA\_AQFT\_Proofs\_v10.lean:

\begin{table}[h]
\centering
\caption{Lean~4 verified results}
\label{tab:lean}
\begin{tabular}{llp{7cm}}
\toprule
ID & Theorem & Statement \\
\midrule
L01 & LLC & $\Sigma_{\mathrm{out}}(v)=\Sigma_{\mathrm{in}}(v)$ at every DAG vertex \\
L02 & Graph Cut & LLC holds independently on each side of causal severance \\
L03 & Markov Blanket & Subgraph shielding from exterior \\
L04 & Frame Independence & $S(U\rho U^\dagger\|\sigma_0)=S(\rho\|\sigma_0)$ for all Poincar\'e $U$ \\
L05 & Rindler Stationarity & $\frac{d}{dt}S(\alpha_t(\rho_\beta)\|\sigma_0)=0$ under modular flow \\
L06 & Rindler Thermal & $\rho_\beta$ is a valid density matrix \\
L07 & Causal Inv.\ (cond.) & $\mu_\sigma(S)=\mu_{\sigma'}(S)$ given amplitude locality \\
L08 & $\alpha_{\mathrm{EM}}$ & $144-7=137$ by \texttt{norm\_num} \\
L09 & Koide & $K=2/3$ from BDG integers \\
L10 & Confinement & $L_{\text{gluon}}=3$, $L_{\text{quark}}=4$ \\
L11 & BDG Closure & 5 topology types, 124 extensions verified \\
L12 & Qubit Fragility & Electrons/photons at minimum BDG score 1 \\
O01 & Amplitude Locality & $a(v|C)$ depends only on $C\cap\mathrm{past}(v)$ \\
O02 & Causal Inv.\ (full) & Invariant for any permutation $\sigma.\mathrm{Perm}\;\sigma'$ \\
O10s & Spin-2 DOF & $5-1-1-1=2$ physical degrees of freedom \\
\bottomrule
\end{tabular}
\end{table}

One intentional sorry remains: a Lean-QuantumInfo library adapter in RA\_AQFT\_Proofs\_v10.lean, unrelated to any RA claim.  Total verified results: $\approx110$.

%───────────────────────────────────────────────────────────────
\section{The Coherence Hub}\label{sec:hub}

\subsection{The Universal Timescale}

The universal formula $\tact = \DeltaS/(\Gamma_{\mathrm{on\text{-}shell}}\cdot\Delta S_{\mathrm{per\;boson}})$ is the single expression connecting all scales:
\begin{align}
  T=1\,\mathrm{mK}: &\quad \tact \approx 1.9\times10^4\,\mathrm{s} \gg \tau_{\mathrm{exp}} &&\text{(quantum regime — BMV null)} \\
  T=300\,\mathrm{K}: &\quad \tau_d \approx 20\,\mathrm{fs} = \tau_{\mathrm{vib}} &&\text{(classical regime — biology)} \\
  \text{Rindler}: &\quad \tact \to \infty &&\text{(KMS: $\Delta S_{\mathrm{per\;photon}}=0$ --- Unruh resolution)} \\
  N=2\text{ seed}: &\quad \tau_d = 0 &&\text{(instantaneous bootstrap)} \\
  \text{Spin-bath}: &\quad t^* = 0.274/g &&\text{(parameter-free prediction)}
\end{align}
Fourteen orders of magnitude.  No free parameters.

\subsection{The $N_{\min}=2$ Identity}

The minimum viable universe seed (Paper~III) and the Causal Firewall origin-of-life threshold (Paper~IV) are the same condition: $N\geq 2$ causally connected actualization events with a directed edge.  The Big Bang and the origin of life share the same physical threshold.

\subsection{The D4U02 Hub}

The D4U02 analytic proof \statusDR\ (Paper~I, Theorem~1) established that the modular barrier $9\equiv 1\pmod{8}$ is specific to $d=4$.  This single result unlocks five others: (1) dimensional uniqueness from the conditioning method; (2) O05 (Unruh detector, via $\tau_{\mathrm{act}}$); (3) O08 (Firewall $\tau_d$, same formula at molecular scale); (4) new perturbative path for O03; (5) the modular invariants question for Lorentz emergence.

%───────────────────────────────────────────────────────────────
\section{Quantum Information}

\subsection{Kinematic Coherence Bound}

\begin{leanbox}
\textbf{Theorem (L12) \statusLV:} Electrons and photons have minimum BDG stability score $1$ (\texttt{structural\_fragility}, RA\_D1\_Proofs.lean).
\end{leanbox}

Each qubit is a kinematic site for spurious actualization.  The structural ceiling on fault-tolerant quantum arrays:
\begin{equation}
  N_{\max} = \eta \cdot p_{\mathrm{th}}
  \label{eq:kcb}
\end{equation}
where $\eta$ is the single-qubit quality factor and $p_{\mathrm{th}}$ the fault-tolerance threshold.  Beyond $N_{\max}$, the cumulative probability of spurious actualization exceeds the fault-tolerance budget regardless of single-qubit fidelity.  Standard QM predicts exponentially difficult but unbounded scaling; RA predicts a hard ceiling.  Near-future QEC experiments at $10^2$--$10^3$ logical qubits will test this.

%───────────────────────────────────────────────────────────────
\section{Open Problems and New Avenues}

\begin{table}[h]
\centering
\caption{Open problems with proof paths}
\label{tab:open}
\begin{tabular}{lp{5cm}p{6cm}}
\toprule
ID & Problem & New avenue (2026) \\
\midrule
O03 & Bianchi on curved background & Perturbative $(l_P/R)^2$ expansion; curvature correction bounded by LLC range (L01+L02) — no Tomita-Takesaki needed at leading order \\
O06 & $l_{\mathrm{RA}}$ from first principles & Dominant $k=2$ term in D4U02 corresponds to minimal BDG diamond: area $= l_P^2 \Rightarrow l_{\mathrm{RA}}=l_P$ \\
O10-G & Lean: $G_{\mu\nu}$ derivation & RA-native route: D04 (Bianchi flat) + O10-spin2 (LV) directly, bypassing Weinberg's S-matrix theorem \\
O11 & Lorentz invariance from DAG & BDG modular invariants ($c_k \bmod \gcd$) as discrete Lorentz-invariants; new question for Dowker collaboration \\
\bottomrule
\end{tabular}
\end{table}

%───────────────────────────────────────────────────────────────
\section{All Predictions}

\begin{table}[h]
\centering
\caption{RA predictions with experimental status}
\label{tab:predictions}
\begin{tabular}{lp{5cm}ll}
\toprule
ID & Prediction & Status & Experiment \\
\midrule
P01 & BMV null; $\tact \gg \tau_{\mathrm{exp}}$ at $1\,\mathrm{mK}$ & \statusDR & BMV groups, near-term \\
P02 & $H_0 \propto \rho_b$(LOS); $75.9$ at Eridanus & \statusCV & DESI, current \\
P03 & No DM below neutrino floor & \statusAR & LZ/XENONnT, ongoing \\
P04 & $t^*=0.274/g$ (spin-bath) & \statusCV & SC circuits \\
P05 & $N_{\max}=\eta p_{\mathrm{th}}$ (QEC ceiling) & \statusLV & QEC, medium-term \\
P06 & Life wherever F1+F2; substrate-independent & \statusAR & SETI, long-term \\
P07 & $\Sigma m_\nu \approx 59\,\mathrm{meV}$, normal hierarchy & conj. & CMB-S4/Euclid \\
P08 & No proton decay (baryon number topological) & \statusDR & — \\
P09 & $\theta_{\mathrm{QCD}}=0$ exactly (DAG acyclic) & \statusDR & — \\
\bottomrule
\end{tabular}
\end{table}

%───────────────────────────────────────────────────────────────
\section{The Laws are the Phenomena}

Physics has always assumed a hierarchy: laws govern phenomena.  RA dissolves this.  The Local Ledger Condition is not a constraint that actualization events satisfy; it \emph{is} what actualization events are, stated as a bookkeeping property.  Conservation of energy is not imposed from above; it is the LLC read at the appropriate scale.  The arrow of time is not explained by initial conditions; it is encoded in the definition of actualization.

The open problems (Table~\ref{tab:open}) are not independent mysteries.  Each is the same question --- ``when does a process become irreversible?'' --- asked at a different scale.  The answer is always $\DeltaS = 0.60069\,\mathrm{nats}$, because that question has only one answer when irreversibility is the primitive.

\section*{Acknowledgments}
This programme was developed over twenty-three years and formalized over approximately twelve months with AI collaboration (Claude, Anthropic; Gemini, Google).  The Lean~4 infrastructure references the Lean-QuantumInfo library (Meiburg/Perimeter).  All papers open-access at \url{https://zenodo.org/communities/relational-actualism}.

\bibliographystyle{plain}
\begin{thebibliography}{99}
\bibitem{RA_P1} Sandeman, J.~F. (2026). Paper~I: Coupling Constants. Zenodo, doi:10.5281/zenodo.19362289.
\bibitem{RA_P2} Sandeman, J.~F. (2026). Paper~II: Wave Function Collapse. Zenodo, doi:10.5281/zenodo.19362968.
\bibitem{RA_P3} Sandeman, J.~F. (2026). Paper~III: Cosmology. Zenodo, doi:10.5281/zenodo.19363017.
\bibitem{RA_P4} Sandeman, J.~F. (2026). Paper~IV: The Causal Firewall. Zenodo, doi:10.5281/zenodo.19198224.
\bibitem{Dowker2013} Dowker, F., and Glaser, L. (2013). Causal set d'Alembertians for various dimensions. \textit{CQG}, 30, 195016.
\bibitem{Weinberg1964} Weinberg, S. (1964). Photons and gravitons in S-matrix theory. \textit{Phys.\ Rev.}, 135, B1049.
\bibitem{Bombelli2009} Bombelli, L., Henson, J., and Sorkin, R.~D. (2009). Discreteness without symmetry breaking. \textit{Mod.\ Phys.\ Lett.\ A}, 24, 2579.
\bibitem{Smolin1992} Smolin, L. (1992). Did the universe evolve? \textit{CQG}, 9, 173.
\bibitem{Mathlib} Mathlib Contributors (2024). Lean~4 Mathematical Library. \url{https://leanprover-community.github.io/mathlib4_docs/}
\end{thebibliography}
\end{document}
